\chapter{Dénombrer sans compter deux fois}
\section*{Exercice 3.1 -- Codes}
[Base] Compter les codes de quatre chiffres distincts, zéros initiaux autorisés. Puis les nombres de quatre chiffres distincts.

\section*{Exercice 3.2 -- Comité}
[Base] Parmi 10 personnes, former un comité de 4 dont deux personnes fixées doivent faire partie.

\section*{Exercice 3.3 -- Au moins un}
[Méthode] Parmi 5 femmes et 4 hommes, combien de groupes de 3 comportent au moins une femme ?

\section*{Exercice 3.4 -- Réunion}
[Base] Sur 40 élèves, 22 font A, 18 font B et 9 font les deux. Compter ceux qui font exactement une activité, puis aucune.

\section*{Exercice 3.5 -- Étoiles et séparateurs}
[Approfondissement] Compter les triplets d'entiers positifs ou nuls de somme 6.

\section*{Exercice 3.6 -- Sous-ensembles}
[Preuve] Dans un ensemble de $n\ge1$ éléments, combien de parties contiennent un élément fixé ?
